% ===========================================================
\chapter{Clase 4: Funciones, Equipotencia e Infinitud — Primeros Resultados de Cantor}
% ===========================================================

La noción de \emph{cardinalidad} —el «tamaño» de un conjunto— adquiere su forma definitiva cuando se formaliza mediante \textbf{biyecciones}. La idea central, debida a Georg Cantor en la década de 1870, es radical en su simplicidad: dos conjuntos tienen el mismo número de elementos si y solo si existe una correspondencia uno a uno entre ellos, independientemente de que los conjuntos sean finitos o infinitos. Esta sesión formaliza dicha idea, estudia sus consecuencias inmediatas —la equipotencia como relación de equivalencia y el Teorema de Cantor–Bernstein–Schröder— y culmina con el descubrimiento más sorprendente de Cantor: existen \emph{distintos tamaños de infinito}, en el sentido de que $\mathbb{R}$ no puede ponerse en correspondencia biyectiva con $\mathbb{N}$.

Todo el desarrollo sigue la tradición de la teoría de conjuntos axiomática ZFC tal como se estableció en la Clase~2. En particular, una función es, por definición, un conjunto de pares ordenados que satisface ciertas condiciones \cite{enderton1977, jech2003, kunen2011}.

% ===========================================================
\section{Funciones: definición formal y tipos}
% ===========================================================

\subsection{Funciones como relaciones}

\begin{definicion}{Función}{funcion}
Una \textbf{función} (o \textbf{aplicación}) de $A$ en $B$ es una relación $f \subseteq A \times B$ tal que:
\begin{enumerate}[leftmargin=2em, label=(\roman*)]
  \item \textbf{Totalidad:} $\forall a \in A,\; \exists b \in B,\; (a,b) \in f$.
  \item \textbf{Unicidad:} $\forall a \in A,\; \forall b, b' \in B,\; \bigl[(a,b) \in f \wedge (a,b') \in f \to b = b'\bigr]$.
\end{enumerate}
Se escribe $f: A \to B$ y, cuando $(a,b) \in f$, se escribe $f(a) = b$. El conjunto $A$ se llama el \textbf{dominio} de $f$ y $B$ el \textbf{codominio}. La \textbf{imagen} (o \textbf{rango}) de $f$ es
\[
  \mathrm{Im}(f) \;=\; f(A) \;=\; \{f(a) \mid a \in A\} \;=\; \{b \in B \mid \exists a \in A,\, f(a) = b\}.
\]
\end{definicion}

\begin{observacion}{Las funciones son conjuntos}{funcconjunto}
En ZFC, toda función $f: A \to B$ es literalmente un conjunto: el conjunto de pares de Kuratowski $\{(a, f(a)) \mid a \in A\}$. La totalidad y la unicidad garantizan que este conjunto «asigna» exactamente un elemento de $B$ a cada elemento de $A$. El conjunto de todas las funciones de $A$ en $B$ se denota $B^A$ o ${}^{A}B$, y su existencia se garantiza mediante el Axioma de Potencia y el Esquema de Separación (pues $B^A \subseteq \mathcal{P}(A \times B)$).
\end{observacion}

\begin{definicion}{Imagen directa e imagen inversa}{imagenes}
Sea $f: A \to B$. Para $S \subseteq A$ y $T \subseteq B$, se definen:
\begin{itemize}[leftmargin=2em]
  \item \textbf{Imagen directa:} $f(S) = \{f(a) \mid a \in S\} \subseteq B$.
  \item \textbf{Preimagen} (o \textbf{imagen inversa}): $f^{-1}(T) = \{a \in A \mid f(a) \in T\} \subseteq A$.
\end{itemize}
\end{definicion}

\begin{ejemplo}{Imágenes directas e inversas}{}
Sea $f: \mathbb{Z} \to \mathbb{Z}$ definida por $f(n) = n^2$.
\begin{enumerate}[leftmargin=2em, label=(\alph*)]
  \item $f(\{-2, -1, 0, 1, 2\}) = \{0, 1, 4\}$.
  \item $f(\mathbb{Z}) = \{0, 1, 4, 9, 16, \ldots\} = \{n^2 \mid n \in \mathbb{Z}\}$, el conjunto de cuadrados perfectos no negativos.
  \item $f^{-1}(\{4\}) = \{-2, 2\}$.
  \item $f^{-1}(\{3\}) = \emptyset$ (pues $3$ no es cuadrado perfecto).
  \item $f^{-1}(\{0, 1, 4\}) = \{-2, -1, 0, 1, 2\}$.
\end{enumerate}
Nótese que $f^{-1}(\{4\})$ tiene dos elementos: $f$ no es inyectiva.
\end{ejemplo}

\begin{proposicion}{Propiedades de las imágenes}{propimag}
Sea $f: A \to B$, $S, S' \subseteq A$ y $T, T' \subseteq B$. Entonces:
\begin{enumerate}[leftmargin=2em, label=(\roman*)]
  \item $f(S \cup S') = f(S) \cup f(S')$.
  \item $f(S \cap S') \subseteq f(S) \cap f(S')$ \quad (la igualdad no siempre se cumple).
  \item $f^{-1}(T \cup T') = f^{-1}(T) \cup f^{-1}(T')$.
  \item $f^{-1}(T \cap T') = f^{-1}(T) \cap f^{-1}(T')$.
  \item $f^{-1}(B \setminus T) = A \setminus f^{-1}(T)$.
  \item $S \subseteq f^{-1}(f(S))$; la igualdad se cumple si $f$ es inyectiva.
  \item $f(f^{-1}(T)) \subseteq T$; la igualdad se cumple si $f$ es sobreyectiva.
\end{enumerate}
\end{proposicion}

\begin{proof}
Demostramos (i) y (ii); las demás son análogas.

\textbf{(i)} Sea $b \in f(S \cup S')$. Entonces existe $a \in S \cup S'$ con $f(a) = b$. Si $a \in S$, entonces $b \in f(S)$; si $a \in S'$, entonces $b \in f(S')$. En cualquier caso, $b \in f(S) \cup f(S')$. La otra inclusión es evidente.

\textbf{(ii)} Si $b \in f(S \cap S')$, existe $a \in S \cap S'$ con $f(a) = b$; entonces $a \in S$ y $a \in S'$, luego $b \in f(S)$ y $b \in f(S')$. Para ver que la inclusión puede ser estricta, tomemos $A = \{1, 2\}$, $B = \{0\}$, $f(1) = f(2) = 0$, $S = \{1\}$, $S' = \{2\}$. Entonces $f(S \cap S') = f(\emptyset) = \emptyset$ pero $f(S) \cap f(S') = \{0\} \cap \{0\} = \{0\}$.
\end{proof}

\subsection{Inyectividad, sobreyectividad y biyectividad}

\begin{definicion}{Función inyectiva}{inyectiva}
Una función $f: A \to B$ es \textbf{inyectiva} (o \textbf{uno a uno}) si
\[
  \forall a, a' \in A,\quad f(a) = f(a') \;\implies\; a = a'.
\]
Equivalentemente, $a \neq a'$ implica $f(a) \neq f(a')$.
\end{definicion}

\begin{definicion}{Función sobreyectiva}{sobreyectiva}
Una función $f: A \to B$ es \textbf{sobreyectiva} (o \textbf{suprayectiva}, o \emph{sobre}) si
\[
  \forall b \in B,\quad \exists a \in A,\quad f(a) = b,
\]
es decir, $\mathrm{Im}(f) = B$.
\end{definicion}

\begin{definicion}{Función biyectiva}{biyectiva}
Una función $f: A \to B$ es \textbf{biyectiva} (o una \textbf{biyección}, o \textbf{correspondencia uno a uno}) si es simultáneamente inyectiva y sobreyectiva.
\end{definicion}

\begin{observacion}{Intuición geométrica}{intuicionbiy}
\begin{itemize}[leftmargin=2em]
  \item \emph{Inyectiva:} distintas entradas producen distintas salidas; «no colapsa elementos».
  \item \emph{Sobreyectiva:} todo elemento del codominio es «alcanzado»; ningún elemento de $B$ queda sin preimagen.
  \item \emph{Biyectiva:} empareja perfectamente cada elemento de $A$ con exactamente uno de $B$ y viceversa; es el concepto preciso de «misma cantidad de elementos».
\end{itemize}
\end{observacion}

\begin{ejemplo}{Clasificación de funciones}{}
\begin{enumerate}[leftmargin=2em, label=(\alph*)]
  \item $f: \mathbb{N} \to \mathbb{N}$, $f(n) = 2n$. Es \textbf{inyectiva} (si $2n = 2m$ entonces $n = m$) pero \textbf{no sobreyectiva} ($1 \notin \mathrm{Im}(f)$).

  \item $g: \mathbb{Z} \to \mathbb{Z}$, $g(n) = n^2$. Es \textbf{no inyectiva} ($g(-2) = g(2) = 4$) y \textbf{no sobreyectiva} ($-1 \notin \mathrm{Im}(g)$).

  \item $h: \mathbb{R} \to \mathbb{R}^+$, $h(x) = e^x$. Es \textbf{inyectiva} (la exponencial es estrictamente creciente) y \textbf{sobreyectiva} (todo $y > 0$ tiene un logaritmo), luego es \textbf{biyectiva}.

  \item La función \textbf{identidad} $\mathrm{id}_A: A \to A$, $\mathrm{id}_A(a) = a$ es siempre biyectiva.

  \item Sea $A = \{1, 2, 3\}$, $B = \{a, b, c\}$. La función $f(1) = a$, $f(2) = a$, $f(3) = b$ es \textbf{no inyectiva} y \textbf{no sobreyectiva} ($c \notin \mathrm{Im}(f)$).

  \item La función \textbf{de Cantor} $\langle \cdot, \cdot \rangle: \mathbb{N} \times \mathbb{N} \to \mathbb{N}$, dada por
  \[
    \langle m, n \rangle = \frac{(m+n)(m+n+1)}{2} + m,
  \]
  es biyectiva. Esto demuestra que $|\mathbb{N} \times \mathbb{N}| = |\mathbb{N}|$, un hecho aparentemente paradójico.
\end{enumerate}
\end{ejemplo}

\begin{ejemplo}{La función de paridad}{paridad}
Definamos $f: \mathbb{N} \to \{0,1\}$ por $f(n) = n \bmod 2$ (el resto al dividir por 2). Entonces:
\begin{itemize}[leftmargin=2em]
  \item $f$ es \textbf{sobreyectiva}: $f(0) = 0$ y $f(1) = 1$.
  \item $f$ \textbf{no} es inyectiva: $f(0) = f(2) = 0$.
\end{itemize}
Las preimágenes son los conjuntos de pares e impares: $f^{-1}(\{0\}) = \{0, 2, 4, \ldots\}$ y $f^{-1}(\{1\}) = \{1, 3, 5, \ldots\}$.
\end{ejemplo}

\subsection{Composición e inversas}

\begin{definicion}{Composición de funciones}{composicion}
Sean $f: A \to B$ y $g: B \to C$. La \textbf{composición} $g \circ f: A \to C$ se define por
\[
  (g \circ f)(a) = g(f(a)) \qquad \forall a \in A.
\]
Como conjunto: $g \circ f = \{(a, c) \in A \times C \mid \exists b \in B,\, (a,b) \in f \wedge (b,c) \in g\}$.
\end{definicion}

\begin{proposicion}{Composición y tipo de función}{compoinysupr}
Sea $f: A \to B$ y $g: B \to C$.
\begin{enumerate}[leftmargin=2em, label=(\roman*)]
  \item Si $f$ y $g$ son inyectivas, entonces $g \circ f$ es inyectiva.
  \item Si $f$ y $g$ son sobreyectivas, entonces $g \circ f$ es sobreyectiva.
  \item Si $f$ y $g$ son biyectivas, entonces $g \circ f$ es biyectiva.
  \item Si $g \circ f$ es inyectiva, entonces $f$ es inyectiva.
  \item Si $g \circ f$ es sobreyectiva, entonces $g$ es sobreyectiva.
\end{enumerate}
\end{proposicion}

\begin{proof}
\textbf{(i)} Si $(g \circ f)(a) = (g \circ f)(a')$, entonces $g(f(a)) = g(f(a'))$. Como $g$ es inyectiva, $f(a) = f(a')$. Como $f$ es inyectiva, $a = a'$.

\textbf{(ii)} Sea $c \in C$. Como $g$ es sobreyectiva, $\exists b \in B$ con $g(b) = c$. Como $f$ es sobreyectiva, $\exists a \in A$ con $f(a) = b$. Entonces $(g \circ f)(a) = g(b) = c$.

\textbf{(iii)} Sigue de (i) y (ii).

\textbf{(iv)} Si $f(a) = f(a')$, entonces $g(f(a)) = g(f(a'))$, es decir, $(g \circ f)(a) = (g \circ f)(a')$. Como $g \circ f$ es inyectiva, $a = a'$.

\textbf{(v)} Sea $c \in C$. Como $g \circ f$ es sobreyectiva, $\exists a \in A$ con $(g \circ f)(a) = c$, es decir, $g(f(a)) = c$. Tomando $b = f(a) \in B$, tenemos $g(b) = c$.
\end{proof}

\begin{definicion}{Función inversa}{inversa}
Sea $f: A \to B$ biyectiva. La \textbf{función inversa} $f^{-1}: B \to A$ se define por
\[
  f^{-1}(b) = a \quad \iff \quad f(a) = b.
\]
La unicidad de $a$ está garantizada por la inyectividad de $f$; la existencia, por la sobreyectividad.
\end{definicion}

\begin{proposicion}{Propiedades de la inversa}{propinversa}
Sea $f: A \to B$ biyectiva. Entonces:
\begin{enumerate}[leftmargin=2em, label=(\roman*)]
  \item $f^{-1}$ es biyectiva y $(f^{-1})^{-1} = f$.
  \item $f^{-1} \circ f = \mathrm{id}_A$ y $f \circ f^{-1} = \mathrm{id}_B$.
  \item Si $g: B \to C$ también es biyectiva, entonces $(g \circ f)^{-1} = f^{-1} \circ g^{-1}$.
\end{enumerate}
\end{proposicion}

\begin{proof}
\textbf{(i)} $f^{-1}$ es inyectiva: si $f^{-1}(b) = f^{-1}(b')$, sea $a$ ese valor común; entonces $f(a) = b$ y $f(a) = b'$, luego $b = b'$. $f^{-1}$ es sobreyectiva: para todo $a \in A$, $f(a) \in B$ y $f^{-1}(f(a)) = a$.

\textbf{(ii)} $(f^{-1} \circ f)(a) = f^{-1}(f(a)) = a$ por definición de $f^{-1}$. Análogamente $(f \circ f^{-1})(b) = f(f^{-1}(b)) = b$.

\textbf{(iii)} $(g \circ f)(a) = c \iff g(f(a)) = c \iff f(a) = g^{-1}(c) \iff a = f^{-1}(g^{-1}(c))$.
\end{proof}

% ===========================================================
\section{Conjuntos finitos e infinitos}
% ===========================================================

\subsection{Conjuntos finitos y la numeración por naturales}

\begin{definicion}{Segmentos iniciales de $\mathbb{N}$}{segnaturales}
Para cada $n \in \mathbb{N}$, definimos el \textbf{segmento inicial}
\[
  [n] \;=\; \{k \in \mathbb{N} \mid k < n\} \;=\; \{0, 1, 2, \ldots, n-1\}.
\]
En particular, $[0] = \emptyset$ y $[n]$ tiene exactamente $n$ elementos.
\end{definicion}

\begin{definicion}{Conjunto finito}{conjfinito}
Un conjunto $A$ es \textbf{finito} si existe $n \in \mathbb{N}$ y una biyección $f: [n] \to A$. En ese caso, el \textbf{cardinal} de $A$ es $|A| = n$. Un conjunto que no es finito se dice \textbf{infinito}.
\end{definicion}

\begin{observacion}{El cardinal está bien definido}{cardinalbien}
Para que la Definición~\ref{def:conjfinito} sea coherente, hay que verificar que si existen biyecciones $f: [m] \to A$ y $g: [n] \to A$, entonces $m = n$. Esto es el llamado \textbf{Principio de las Casillas} (o lema de Dedekind–Cantor para conjuntos finitos) y se prueba por inducción sobre $\min(m,n)$. No lo demostraremos en detalle aquí, pero el resultado es esencial para la consistencia de la noción de cardinalidad finita.
\end{observacion}

\begin{proposicion}{Propiedades de los conjuntos finitos}{propfinitos}
\begin{enumerate}[leftmargin=2em, label=(\roman*)]
  \item Todo subconjunto de un conjunto finito es finito.
  \item La unión finita de conjuntos finitos es finita. En particular, si $|A| = m$ y $|B| = n$ y $A \cap B = \emptyset$, entonces $|A \cup B| = m + n$.
  \item Si $|A| = m$ y $|B| = n$, entonces $|A \times B| = m \cdot n$.
  \item Si $|A| = n$, entonces $|\mathcal{P}(A)| = 2^n$.
  \item No existe biyección entre $[m]$ y $[n]$ para $m \neq n$ (principio de las casillas).
\end{enumerate}
\end{proposicion}

\subsection{La definición de Dedekind de conjunto infinito}

La noción de «infinito» puede también caracterizarse desde adentro, sin referencia explícita a $\mathbb{N}$:

\begin{definicion}{Conjunto Dedekind-infinito}{dedinfinito}
Un conjunto $A$ es \textbf{Dedekind-infinito} si existe una biyección $f: A \to B$ donde $B \subsetneq A$ (un subconjunto \textbf{propio} de $A$). Equivalentemente, $A$ es Dedekind-infinito si existe una función inyectiva $f: A \to A$ que \textbf{no} es sobreyectiva.

Un conjunto que no es Dedekind-infinito se llama \textbf{Dedekind-finito}.
\end{definicion}

\begin{observacion}{Origen histórico}{dedekindhistoria}
Esta definición aparece en el trabajo de Richard Dedekind «Was sind und was sollen die Zahlen?» (1888), donde la usa precisamente para \emph{definir} los conjuntos infinitos de manera intrínseca, sin depender de los naturales \cite{dedekind1888}. La idea es que un conjunto es infinito si y solo si tiene «la misma cantidad» de elementos que alguna parte propia suya; algo imposible en los conjuntos finitos.
\end{observacion}

\begin{ejemplo}{$\mathbb{N}$ es Dedekind-infinito}{ndedinfinito}
La función $f: \mathbb{N} \to \mathbb{N}$ dada por $f(n) = n + 1$ es inyectiva (si $n+1 = m+1$ entonces $n = m$) pero no sobreyectiva ($0 \notin \mathrm{Im}(f)$). Su imagen es $\mathbb{N} \setminus \{0\} = \{1, 2, 3, \ldots\} \subsetneq \mathbb{N}$.

Más aún, $f$ establece una biyección entre $\mathbb{N}$ y el subconjunto propio $\{1, 2, 3, \ldots\}$, lo que confirma que $\mathbb{N}$ es Dedekind-infinito. Este fenómeno fue la fuente de la famosa paradoja informal del «Hotel de Hilbert».
\end{ejemplo}

\begin{ejemplo}{El Hotel de Hilbert (ilustración intuitiva)}{hilbert}
Imaginemos un hotel con habitaciones numeradas $0, 1, 2, 3, \ldots$ (una por cada natural), todas ocupadas. Si llega un nuevo huésped:
\begin{itemize}[leftmargin=2em]
  \item Movemos al huésped de la habitación $n$ a la habitación $n+1$ (para todo $n \in \mathbb{N}$).
  \item La habitación $0$ queda libre para el nuevo huésped.
\end{itemize}
Este «movimiento» corresponde exactamente a la función $f(n) = n+1$: es una biyección de $\mathbb{N}$ en $\mathbb{N} \setminus \{0\}$. Formalmente, reacomodamos el hotel estableciendo la correspondencia
\[
  \{0,1,2,3,\ldots\} \;\longleftrightarrow\; \{1,2,3,4,\ldots\},
\]
dejando libre la habitación $0$ para el nuevo cliente. Si llegan infinitos nuevos huéspedes (numerados $0, 1, 2, \ldots$), movemos al huésped de la habitación $n$ a la $2n+1$ (impar), dejando todas las pares libres. Esta es la idea detrás de muchas pruebas de numerabilidad.
\end{ejemplo}

\begin{proposicion}{Conjuntos finitos son Dedekind-finitos}{finDedfinito}
Todo conjunto finito es Dedekind-finito; es decir, ningún conjunto finito es biyectable con un subconjunto propio suyo.
\end{proposicion}

\begin{proof}
Por inducción en $n = |A|$. Si $|A| = 0$, entonces $A = \emptyset$ y no tiene subconjuntos propios. Supongamos el resultado cierto para todos los conjuntos de cardinal menor que $n$, y sea $|A| = n \geq 1$. Supongamos que $f: A \to B$ es una biyección con $B \subsetneq A$. Sea $a_0 \in A \setminus B$. Definimos $A' = A \setminus \{a_0\}$, que tiene $n-1$ elementos. La restricción de $f$ da una inyección de $A$ en $A'$ (pues $B \subseteq A' \cup B \subseteq A$). En particular, $|f(A)| = |A| = n$ pero $f(A) = B \subseteq A' \setminus \{a_0\}$, que tiene a lo sumo $n-1$ elementos. Esto contradice el principio de las casillas. (Una prueba detallada requiere el principio de las casillas, que a su vez se demuestra por inducción.)
\end{proof}

\begin{proposicion}{Relación entre infinitud e infinitud de Dedekind}{infDed}
Bajo el Axioma de Elección (AC), las siguientes condiciones sobre un conjunto $A$ son equivalentes:
\begin{enumerate}[leftmargin=2em, label=(\roman*)]
  \item $A$ es infinito (no finito).
  \item $A$ es Dedekind-infinito.
  \item Existe una función inyectiva $\mathbb{N} \to A$.
\end{enumerate}
Sin el Axioma de Elección, puede existir un conjunto infinito que sea Dedekind-finito. \cite{herrlich2006}
\end{proposicion}

\begin{proof}[Esquema de la prueba]
\textbf{(iii) $\Rightarrow$ (ii):} Si $\iota: \mathbb{N} \to A$ es inyectiva, definimos $f: A \to A$ por
\[
  f(a) = \begin{cases} \iota(\iota^{-1}(a) + 1) & \text{si } a \in \mathrm{Im}(\iota), \\ a & \text{si } a \notin \mathrm{Im}(\iota). \end{cases}
\]
Entonces $f$ es inyectiva y $\iota(0) \notin \mathrm{Im}(f)$, luego $A$ es Dedekind-infinito.

\textbf{(ii) $\Rightarrow$ (iii):} Si $f: A \to A$ es inyectiva y $a_0 \in A \setminus \mathrm{Im}(f)$, definimos $\iota: \mathbb{N} \to A$ por $\iota(0) = a_0$ e $\iota(n+1) = f(\iota(n))$. Por inducción, $\iota$ es inyectiva.

\textbf{(i) $\Rightarrow$ (iii):} Usando AC (de forma controlada), un conjunto infinito $A$ contiene subconjuntos de todo tamaño finito, y se puede construir por inducción una función inyectiva de $\mathbb{N}$ en $A$.

\textbf{(iii) $\Rightarrow$ (i):} Si existiera una biyección $h: [n] \to A$ y una inyección $\iota: \mathbb{N} \to A$, la composición $h^{-1} \circ \iota: \mathbb{N} \to [n]$ sería inyectiva, lo que contradice el principio de las casillas.
\end{proof}

% ===========================================================
\section{Equipotencia como relación de equivalencia}
% ===========================================================

\subsection{Definición y primeras propiedades}

\begin{definicion}{Equipotencia}{equipotencia}
Dos conjuntos $A$ y $B$ son \textbf{equipotentes} (o tienen el mismo \textbf{cardinal}) si existe una biyección $f: A \to B$. Se escribe $A \sim B$ o $|A| = |B|$.
\end{definicion}

\begin{teorema}{La equipotencia es una relación de equivalencia}{equivequipot}
La relación $\sim$ de equipotencia satisface:
\begin{enumerate}[leftmargin=2em, label=(\roman*)]
  \item \textbf{Reflexividad:} $A \sim A$ para todo conjunto $A$.
  \item \textbf{Simetría:} Si $A \sim B$, entonces $B \sim A$.
  \item \textbf{Transitividad:} Si $A \sim B$ y $B \sim C$, entonces $A \sim C$.
\end{enumerate}
\end{teorema}

\begin{proof}
\textbf{(i)} La identidad $\mathrm{id}_A: A \to A$ es una biyección.

\textbf{(ii)} Si $f: A \to B$ es biyectiva, entonces $f^{-1}: B \to A$ también lo es (Proposición~\ref{prop:propinversa}).

\textbf{(iii)} Si $f: A \to B$ y $g: B \to C$ son biyectivas, entonces $g \circ f: A \to C$ es biyectiva (Proposición~\ref{prop:compoinysupr}).
\end{proof}

\begin{observacion}{Clases de equivalencia y cardinalidad}{clasecardinal}
Formalmente, la «clase de equivalencia» de $A$ bajo $\sim$ sería el conjunto $\{B \mid B \sim A\}$. Esto no es un conjunto en ZFC (sería demasiado grande, una clase propia). Por ello, la teoría de cardinales introduce los \textbf{números cardinales} como representantes canónicos: en ZFC, el cardinal de $A$ se define como el ordinal mínimo equipotente a $A$ (ordinal de Hartogs). Para los propósitos de esta clase, trabajaremos con la equipotencia directamente sin necesidad de este refinamiento.
\end{observacion}

\begin{ejemplo}{Biyecciones entre conjuntos infinitos}{}
\begin{enumerate}[leftmargin=2em, label=(\alph*)]
  \item $\mathbb{N} \sim \mathbb{Z}$: la función $f: \mathbb{N} \to \mathbb{Z}$ definida por
  \[
    f(n) = \begin{cases} n/2 & \text{si } n \text{ es par}, \\ -(n+1)/2 & \text{si } n \text{ es impar}, \end{cases}
  \]
  es una biyección. Explícitamente: $f(0) = 0$, $f(1) = -1$, $f(2) = 1$, $f(3) = -2$, $f(4) = 2$, $\ldots$

  \item $\mathbb{N} \sim \mathbb{N} \times \mathbb{N}$: la función de emparejamiento de Cantor $\langle m, n \rangle = \frac{(m+n)(m+n+1)}{2} + m$ es biyectiva. Los primeros valores son:
  \[
    \langle 0,0\rangle = 0,\quad \langle 1,0\rangle = 1,\quad \langle 0,1\rangle = 2,\quad \langle 2,0\rangle = 3,\quad \langle 1,1\rangle = 4,\quad \langle 0,2\rangle = 5,\ldots
  \]

  \item $(0,1) \sim \mathbb{R}$: la función $f: (0,1) \to \mathbb{R}$, $f(x) = \tan\!\bigl(\pi(x - \tfrac{1}{2})\bigr)$, es una biyección.

  \item $[0,1] \sim (0,1)$: puede construirse una biyección aunque los dos intervalos no sean idénticos. Una construcción explícita es:
  \[
    f(x) = \begin{cases}
      1/(n+2) & \text{si } x = 1/n \text{ para algún } n \in \mathbb{N}^+, \\
      0 & \text{si } x = 0, \\
      x & \text{en otro caso.}
    \end{cases}
  \]
  Verificar que $f: [0,1] \to (0,1)$ es biyectiva es un buen ejercicio.
\end{enumerate}
\end{ejemplo}

\subsection{Comparación de cardinales: la relación $\leq$}

\begin{definicion}{Comparación de cardinales}{cardleq}
Se dice que $|A| \leq |B|$ (el cardinal de $A$ es \textbf{menor o igual} que el de $B$) si existe una función inyectiva $f: A \to B$. Se escribe $A \preceq B$.

Se dice que $|A| < |B|$ (el cardinal de $A$ es \textbf{estrictamente menor}) si $A \preceq B$ pero $A \not\sim B$.
\end{definicion}

\begin{observacion}{Reflexividad y transitividad de $\preceq$}{reflextransprec}
La relación $\preceq$ es reflexiva (la identidad es inyectiva) y transitiva (la composición de inyecciones es inyectiva). La antisimetría —que $A \preceq B$ y $B \preceq A$ impliquen $A \sim B$— es precisamente el contenido del Teorema de Cantor–Bernstein–Schröder que demostraremos en la siguiente sección.
\end{observacion}

% ===========================================================
\section{Teorema de Cantor–Bernstein–Schröder}
% ===========================================================

\subsection{Enunciado}

\begin{teorema}{Cantor–Bernstein–Schröder}{CBS}
Sean $A$ y $B$ conjuntos. Si existe una inyección $f: A \to B$ y una inyección $g: B \to A$, entonces existe una biyección $h: A \to B$. En notación de cardinales:
\[
  |A| \leq |B| \wedge |B| \leq |A| \;\implies\; |A| = |B|.
\]
\end{teorema}

\begin{observacion}{Historia del teorema}{historiaCBS}
Este resultado fue enunciado por Cantor pero no demostrado satisfactoriamente por él. La primera demostración correcta es debida a Felix Bernstein (1898, en su tesis doctoral) y a Ernst Schröder (1896, aunque con un error subsanado). Existen múltiples demostraciones; presentamos la clásica debida a la idea de los «cadenas de Bernstein» \cite{jech2003, halmos1960, hrbacek1999}.
\end{observacion}

\subsection{Demostración}

\begin{proof}[Demostración del Teorema~\ref{teo:CBS}]
Sean $f: A \to B$ y $g: B \to A$ inyecciones. Construiremos una biyección $h: A \to B$.

\medskip
\textbf{Paso 1: Definición de las cadenas.}

Dado $x \in A \cup B$, definimos la \textbf{cadena} de $x$ como la sucesión (finita o infinita hacia atrás) obtenida aplicando $f^{-1}$ y $g^{-1}$ alternativamente, en tanto sea posible:
\[
  \ldots \to g^{-1}(x) \to f^{-1}(g^{-1}(x)) \to \cdots \to x.
\]
Más precisamente: si $x \in A$, entonces (si existe) $f(x) \in B$, y (si existe) $g^{-1}(x) \in B$, luego (si existe) $f^{-1}(g^{-1}(x)) \in A$, y así sucesivamente. Cada elemento de $A \cup B$ pertenece a exactamente una cadena.

\medskip
\textbf{Paso 2: Tipos de cadenas.}

Cada cadena es de uno de tres tipos:
\begin{itemize}[leftmargin=2em]
  \item \textbf{Tipo A:} la cadena es infinita y el «primer» elemento (el de más a la izquierda) pertenece a $A$.
  \item \textbf{Tipo B:} la cadena es infinita y el primer elemento pertenece a $B$.
  \item \textbf{Tipo $\infty$:} la cadena es doblemente infinita (sin extremo izquierdo).
\end{itemize}

\medskip
\textbf{Paso 3: Definición de $h$.}

Definimos $h: A \to B$ por:
\[
  h(x) = \begin{cases}
    f(x) & \text{si } x \text{ está en una cadena de tipo A o de tipo } \infty, \\
    g^{-1}(x) & \text{si } x \text{ está en una cadena de tipo B.}
  \end{cases}
\]

\medskip
\textbf{Paso 4: $h$ está bien definida.}

Si $x$ está en una cadena de tipo B, entonces $x \in g(\mathrm{Im}(f)) \subseteq \mathrm{Im}(g)$, y por ello $g^{-1}(x)$ existe. En los otros casos, usamos $f(x)$, que siempre está definido.

\medskip
\textbf{Paso 5: $h$ es inyectiva.}

Elementos en cadenas distintas tienen imágenes distintas porque $f$ y $g^{-1}$ son inyecciones. Dentro de la misma cadena, distintos elementos de $A$ producen distintas imágenes por la misma razón.

\medskip
\textbf{Paso 6: $h$ es sobreyectiva.}

Sea $y \in B$. Analizamos el tipo de la cadena de $y$:
\begin{itemize}[leftmargin=2em]
  \item Si $y$ está en una cadena de tipo A o $\infty$: en esa cadena, el predecesor inmediato de $y$ (por $f$) pertenece a $A$; llámalo $x = f^{-1}(y)$. Entonces $h(x) = f(x) = y$.
  \item Si $y$ está en una cadena de tipo B: $g(y) \in A$ está en la misma cadena, y $h(g(y)) = g^{-1}(g(y)) = y$.
\end{itemize}
En todo caso, $y$ tiene preimagen bajo $h$.
\end{proof}

\begin{corolario}{$\preceq$ es antisimétrica}{prec-antisim}
La relación $\preceq$ (de comparación de cardinales por inyecciones) satisface la antisimetría: $A \preceq B$ y $B \preceq A$ implican $A \sim B$.
\end{corolario}

\begin{ejemplo}{Aplicación del Teorema CBS}{aplicCBS}
\begin{enumerate}[leftmargin=2em, label=(\alph*)]
  \item \textbf{$[0,1] \sim (0,1)$:} Tenemos la inyección obvia $\iota: (0,1) \hookrightarrow [0,1]$ (inclusión). Para la otra dirección, la función $f: [0,1] \to (0,1)$ definida por $f(x) = \frac{x+1}{3}$ es inyectiva (pues $\frac{x+1}{3} \in (\frac{1}{3}, \frac{2}{3}) \subset (0,1)$). Por CBS, $[0,1] \sim (0,1)$.

  \item \textbf{$\mathbb{N} \sim 2\mathbb{N}$:} Tenemos $\iota: 2\mathbb{N} \hookrightarrow \mathbb{N}$ (inclusión) y $f: \mathbb{N} \to 2\mathbb{N}$, $f(n) = 2n$ (inyectiva). Por CBS, $\mathbb{N} \sim 2\mathbb{N}$ (también demostrable directamente).

  \item \textbf{$\mathcal{P}(\mathbb{N}) \sim \mathbb{R}$:} Se construyen inyecciones en ambas direcciones. Una dirección: la representación binaria de reales en $(0,1)$ da una inyección de $(0,1)$ (salvo un conjunto numerable) en $\{0,1\}^{\mathbb{N}} \sim \mathcal{P}(\mathbb{N})$; la otra: la representación binaria da una inyección de $\mathcal{P}(\mathbb{N})$ en $[0,1]$. Los detalles se completan con CBS.
\end{enumerate}
\end{ejemplo}

\subsection{Variante constructiva: demostración de Knaster-Tarski}

Existe otra elegante demostración del Teorema CBS que utiliza el Teorema del Punto Fijo de Knaster–Tarski:

\begin{proof}[Demostración alternativa (Knaster--Tarski)]
Sean $f: A \to B$ y $g: B \to A$ inyecciones. Consideremos la función $\Phi: \mathcal{P}(A) \to \mathcal{P}(A)$ definida por
\[
  \Phi(X) \;=\; A \setminus g\!\bigl(B \setminus f(X)\bigr).
\]
Se verifica que $\Phi$ es \textbf{monótona}: si $X \subseteq X'$, entonces $f(X) \subseteq f(X')$, luego $B \setminus f(X) \supseteq B \setminus f(X')$, luego $g(B \setminus f(X)) \supseteq g(B \setminus f(X'))$, luego $\Phi(X) \subseteq \Phi(X')$.

Por el Teorema del Punto Fijo de Knaster–Tarski (el retículo $\mathcal{P}(A)$ es completo), $\Phi$ tiene un \textbf{punto fijo}: existe $C \subseteq A$ con $\Phi(C) = C$, es decir,
\[
  C = A \setminus g\!\bigl(B \setminus f(C)\bigr).
\]
Definimos $h: A \to B$ por
\[
  h(a) = \begin{cases} f(a) & \text{si } a \in C, \\ g^{-1}(a) & \text{si } a \in A \setminus C. \end{cases}
\]
De la ecuación del punto fijo: $A \setminus C = g(B \setminus f(C))$, de modo que $g^{-1}$ está bien definida en $A \setminus C$ con imagen $B \setminus f(C)$. Se comprueba que $h$ es biyectiva: su imagen es $f(C) \cup (B \setminus f(C)) = B$.
\end{proof}

% ===========================================================
\section{Conjuntos numerables y no numerables}
% ===========================================================

\subsection{Conjuntos numerables}

\begin{definicion}{Conjunto numerable}{numerable}
Un conjunto $A$ es:
\begin{itemize}[leftmargin=2em]
  \item \textbf{Finito} si $A \sim [n]$ para algún $n \in \mathbb{N}$.
  \item \textbf{Infinito numerable} (o \textbf{contablemente infinito}) si $A \sim \mathbb{N}$.
  \item \textbf{Numerable} (o \textbf{a lo sumo numerable}) si es finito o infinito numerable, es decir, si $|A| \leq |\mathbb{N}|$.
  \item \textbf{No numerable} (o \textbf{incontable}) si $|\mathbb{N}| < |A|$, es decir, si $A$ no es numerable.
\end{itemize}
\end{definicion}

\begin{observacion}{Terminología alternativa}{terminonumerable}
En la literatura en inglés: \emph{countably infinite} para $A \sim \mathbb{N}$, \emph{countable} para a lo sumo numerable, \emph{uncountable} para no numerable. En español se usa también «contable» como sinónimo de «numerable». El cardinal de $\mathbb{N}$ se denota $\aleph_0$ (aleph-cero), el primero de los \textbf{números cardinales transfinitos} introducidos por Cantor.
\end{observacion}

\begin{teorema}{Caracterizaciones de la numerabilidad}{carnum}
Para un conjunto infinito $A$, las siguientes condiciones son equivalentes:
\begin{enumerate}[leftmargin=2em, label=(\roman*)]
  \item $A \sim \mathbb{N}$.
  \item Existe una enumeración de $A$: una sucesión $a_0, a_1, a_2, \ldots$ tal que $A = \{a_0, a_1, a_2, \ldots\}$ y todos los $a_i$ son distintos.
  \item Existe una función sobreyectiva $\mathbb{N} \to A$.
  \item Existe una función inyectiva $A \to \mathbb{N}$.
\end{enumerate}
\end{teorema}

\begin{proof}
\textbf{(i) $\Leftrightarrow$ (ii):} Una biyección $f: \mathbb{N} \to A$ es exactamente una enumeración $a_n = f(n)$.

\textbf{(i) $\Rightarrow$ (iii):} Toda biyección es sobreyectiva.

\textbf{(iii) $\Rightarrow$ (i):} Sea $s: \mathbb{N} \to A$ sobreyectiva. Definamos $f: A \to \mathbb{N}$ por $f(a) = \min\{n \in \mathbb{N} \mid s(n) = a\}$. Esta $f$ es inyectiva. Como $A$ es infinito, $f(A)$ es un subconjunto infinito de $\mathbb{N}$, y todo subconjunto infinito de $\mathbb{N}$ es equipotente a $\mathbb{N}$ (por la función que enumera sus elementos en orden creciente). Componiendo, $A \sim \mathbb{N}$.

\textbf{(i) $\Rightarrow$ (iv):} Toda biyección $A \to \mathbb{N}$ es inyectiva.

\textbf{(iv) $\Rightarrow$ (i):} Si $\iota: A \hookrightarrow \mathbb{N}$ es inyectiva, entonces $A \sim \iota(A) \subseteq \mathbb{N}$. Como $A$ es infinito y $\iota(A)$ también lo es, $\iota(A)$ es un subconjunto infinito de $\mathbb{N}$, equipotente a $\mathbb{N}$.
\end{proof}

\begin{proposicion}{Todo subconjunto infinito de $\mathbb{N}$ es numerable}{subinfN}
Si $S \subseteq \mathbb{N}$ es infinito, entonces $S \sim \mathbb{N}$.
\end{proposicion}

\begin{proof}
Definimos $f: \mathbb{N} \to S$ por recursión: $f(0) = \min S$ y $f(n+1) = \min(S \setminus \{f(0), \ldots, f(n)\})$. Como $S$ es infinito, este mínimo siempre existe. La función $f$ es estrictamente creciente (y en particular inyectiva) y su imagen es $S$ (si algún $s \in S$ no estuviera en $\mathrm{Im}(f)$, habría una cantidad infinita de elementos de $S$ menores que $s$, lo que es imposible). Luego $f$ es una biyección de $\mathbb{N}$ en $S$.
\end{proof}

\begin{teorema}{Numerabilidad de $\mathbb{Z}$, $\mathbb{Q}$ y uniones numerables}{numZQ}
\begin{enumerate}[leftmargin=2em, label=(\roman*)]
  \item $\mathbb{Z}$ es numerable.
  \item $\mathbb{N} \times \mathbb{N}$ es numerable.
  \item Si $A$ y $B$ son numerables, $A \cup B$ es numerable.
  \item Una unión numerable de conjuntos numerables es numerable: si $(A_n)_{n \in \mathbb{N}}$ es una familia de conjuntos numerables, $\bigcup_{n \in \mathbb{N}} A_n$ es numerable.
  \item $\mathbb{Q}$ es numerable.
\end{enumerate}
\end{teorema}

\begin{proof}
\textbf{(i)} Ya construimos explícitamente la biyección $f: \mathbb{N} \to \mathbb{Z}$.

\textbf{(ii)} La función de emparejamiento de Cantor $\langle m, n \rangle: \mathbb{N} \times \mathbb{N} \to \mathbb{N}$ es biyectiva. Una segunda prueba es por el argumento de la diagonal: enumeramos $\mathbb{N} \times \mathbb{N}$ recorriendo las diagonales $\{(m,n) \mid m + n = k\}$ para $k = 0, 1, 2, \ldots$:
\[
  (0,0),\ (1,0),\ (0,1),\ (2,0),\ (1,1),\ (0,2),\ (3,0),\ \ldots
\]

\textbf{(iii)} Sean $a_0, a_1, a_2, \ldots$ y $b_0, b_1, b_2, \ldots$ enumeraciones de $A$ y $B$ (sin perder generalidad, suponemos $A \cap B = \emptyset$). Entonces $a_0, b_0, a_1, b_1, a_2, b_2, \ldots$ es una enumeración de $A \cup B$.

\textbf{(iv)} Para cada $n$, sea $a_{n,0}, a_{n,1}, a_{n,2}, \ldots$ una enumeración de $A_n$. Definimos una enumeración de $\bigcup A_n$ recorriendo la tabla $(a_{n,k})_{n,k \in \mathbb{N}}$ en diagonal (método de Cantor): ordenamos los pares $(n,k)$ según $n+k$ y luego por $n$, y enumeramos $a_{n,k}$ en ese orden, omitiendo repeticiones. Por (ii), esto da una enumeración.

\textbf{(v)} Tenemos $\mathbb{Q} = \bigcup_{n \in \mathbb{Z}} \{p/n \mid p \in \mathbb{Z}, \gcd(p,n) = 1\}$ (unión numerable de conjuntos numerables). Por (iv), $\mathbb{Q}$ es numerable.

Una demostración más directa: la función $f: \mathbb{Z} \times \mathbb{N}^+ \to \mathbb{Q}$, $f(p,q) = p/q$, es sobreyectiva. Como $\mathbb{Z} \times \mathbb{N}^+$ es numerable (por (ii) y (i)), $\mathbb{Q}$ lo es.
\end{proof}

\begin{ejemplo}{El argumento diagonal para $\mathbb{N} \times \mathbb{N}$}{diagonalNN}
El argumento diagonal de Cantor para $\mathbb{N} \times \mathbb{N}$ puede visualizarse como una tabla infinita:
\[
\begin{array}{cccccc}
(0,0) & (0,1) & (0,2) & (0,3) & (0,4) & \cdots \\
(1,0) & (1,1) & (1,2) & (1,3) & \cdots  & \\
(2,0) & (2,1) & (2,2) & \cdots  &  & \\
(3,0) & (3,1) & \cdots  &  &  & \\
\vdots &  &  &  &  &
\end{array}
\]
Recorremos las antidiagonales $\{(i,j) \mid i+j = k\}$ para $k = 0, 1, 2, \ldots$, obteniendo la enumeración:
\[
(0,0),\ (1,0),\ (0,1),\ (2,0),\ (1,1),\ (0,2),\ (3,0),\ (2,1),\ (1,2),\ (0,3),\ \ldots
\]
\end{ejemplo}

\begin{corolario}{El conjunto de polinomios con coeficientes enteros es numerable}{polnum}
El conjunto $\mathbb{Z}[x]$ de todos los polinomios con coeficientes enteros es numerable.
\end{corolario}

\begin{proof}
Para cada $n \in \mathbb{N}$, el conjunto $\mathbb{Z}_n[x]$ de polinomios de grado $\leq n$ está en biyección con $\mathbb{Z}^{n+1}$ (los coeficientes), que es numerable. Entonces $\mathbb{Z}[x] = \bigcup_{n \geq 0} \mathbb{Z}_n[x]$ es unión numerable de conjuntos numerables, luego numerable.
\end{proof}

\begin{corolario}{Los números algebraicos son numerables}{algnum}
El conjunto $\mathbb{A}$ de los \textbf{números algebraicos} (raíces de polinomios con coeficientes enteros) es numerable.
\end{corolario}

\begin{proof}
Cada polinomio $p \in \mathbb{Z}[x]$ de grado $n$ tiene a lo sumo $n$ raíces complejas. El conjunto total de raíces es $\mathbb{A} = \bigcup_{p \in \mathbb{Z}[x]} \{\text{raíces de } p\}$, unión numerable de conjuntos finitos, luego numerable.
\end{proof}

\begin{observacion}{Números trascendentes}{trascendentes}
Como veremos en la siguiente sección, $\mathbb{R}$ es no numerable. Dado que $\mathbb{A}$ es numerable y $\mathbb{A} \subset \mathbb{R}$, el conjunto $\mathbb{R} \setminus \mathbb{A}$ de los \textbf{números trascendentes} (como $\pi$ y $e$) es también no numerable; en particular, no puede ser vacío. Esto es una demostración de existencia puramente conjuntista: los trascendentes existen (y son «la mayoría» de los reales) sin necesidad de exhibir ningún ejemplo concreto.
\end{observacion}

% ===========================================================
\section{Conjuntos no numerables: el argumento diagonal de Cantor}
% ===========================================================

\subsection{La no numerabilidad de $\mathbb{R}$}

El resultado siguiente es uno de los teoremas más famosos de las matemáticas:

\begin{teorema}{No numerabilidad de $\mathbb{R}$}{nonum}
El conjunto $\mathbb{R}$ de los números reales no es numerable: $|\mathbb{R}| > |\mathbb{N}|$.
\end{teorema}

\begin{proof}[Demostración (argumento diagonal de Cantor, 1891)]
Basta demostrar que el intervalo $(0,1)$ no es numerable (pues $(0,1) \sim \mathbb{R}$ por la biyección $x \mapsto \tan(\pi(x - \frac{1}{2}))$, y si $(0,1)$ fuera numerable también lo sería $\mathbb{R}$).

Supongamos, para obtener una contradicción, que $(0,1)$ \textbf{es} numerable. Entonces existe una enumeración $x_0, x_1, x_2, \ldots$ de todos los elementos de $(0,1)$. Escribimos cada $x_i$ en su expansión decimal:
\[
  x_i = 0.d_{i,0}\, d_{i,1}\, d_{i,2}\, d_{i,3} \cdots
\]
donde $d_{i,k} \in \{0, 1, \ldots, 9\}$ son los dígitos decimales de $x_i$.

Construimos ahora un número $y = 0.e_0\, e_1\, e_2\, e_3\,\cdots \in (0,1)$ definiendo:
\[
  e_k = \begin{cases} 5 & \text{si } d_{k,k} \neq 5, \\ 6 & \text{si } d_{k,k} = 5. \end{cases}
\]
(O cualquier elección que garantice $e_k \neq d_{k,k}$ y evite los dígitos $0$ y $9$ para esquivar ambigüedades en la representación decimal.)

Afirmamos que $y \in (0,1)$ y que $y \neq x_k$ para todo $k \in \mathbb{N}$. En efecto, $y$ y $x_k$ difieren en el $k$-ésimo dígito decimal (en la posición $k$): $e_k \neq d_{k,k}$. Como los dígitos elegidos están en $\{5, 6\}$, no hay ambigüedad de representación (no se presentan los casos $0.999\ldots = 1.000\ldots$), así que esta diferencia de dígito garantiza $y \neq x_k$.

Pero entonces $y \in (0,1)$ y $y \notin \{x_0, x_1, x_2, \ldots\}$, lo cual contradice que la enumeración sea de todos los elementos de $(0,1)$. Por tanto, $(0,1)$ no es numerable.
\end{proof}

\begin{ejemplo}{Construcción explícita de la diagonal}{diagexplicita}
Ilustramos el argumento con una lista hipotética. Supongamos que la enumeración comienza:
\begin{align*}
  x_0 &= 0.\mathbf{3}14159265\ldots \\
  x_1 &= 0.7\mathbf{1}8281828\ldots \\
  x_2 &= 0.57\mathbf{7}21566\ldots \\
  x_3 &= 0.141\mathbf{5}9265\ldots \\
  x_4 &= 0.2718\mathbf{2}8182\ldots \\
  x_5 &= 0.61803\mathbf{3}9887\ldots \\
  \vdots
\end{align*}
Los dígitos de la diagonal son $d_{0,0} = 3$, $d_{1,1} = 1$, $d_{2,2} = 7$, $d_{3,3} = 5$, $d_{4,4} = 2$, $d_{5,5} = 3$, $\ldots$

Aplicando la regla $e_k = 5$ si $d_{k,k} \neq 5$, $e_k = 6$ si $d_{k,k} = 5$:
\[
  y = 0.\underbrace{5}_{e_0}\underbrace{5}_{e_1}\underbrace{5}_{e_2}\underbrace{6}_{e_3}\underbrace{5}_{e_4}\underbrace{5}_{e_5}\cdots = 0.555655\ldots
\]
Este número $y$ difiere de $x_k$ en al menos el $k$-ésimo dígito para cada $k$, luego $y \neq x_k$ para todo $k$.
\end{ejemplo}

\begin{observacion}{Esencia del argumento diagonal}{esenciadiag}
El argumento diagonal es mucho más general que su aplicación a $\mathbb{R}$: es un método para construir, dada cualquier lista de objetos de un cierto tipo, un objeto del mismo tipo que no aparece en la lista. Sus variantes aparecen en:
\begin{itemize}[leftmargin=2em]
  \item El Teorema de Cantor: $|A| < |\mathcal{P}(A)|$ para todo conjunto $A$ (Sección~\ref{sec:cantorpot}).
  \item La prueba de incompletitud de Gödel.
  \item La indecidibilidad del problema de la parada (halting problem) de Turing.
  \item La construcción de números de Liouville (números trascendentes).
\end{itemize}
\end{observacion}

\subsection{Distintos infinitos: $\aleph_0$ y $\mathfrak{c}$}

\begin{definicion}{Los cardinales $\aleph_0$ y $\mathfrak{c}$}{alephnota}
El \textbf{cardinal} $\aleph_0 = |\mathbb{N}|$ es el primer cardinal infinito. El cardinal $\mathfrak{c} = |\mathbb{R}|$ se llama la \textbf{potencia del continuo}. El Teorema~\ref{teo:nonum} establece que $\aleph_0 < \mathfrak{c}$.
\end{definicion}

\begin{proposicion}{Valor de $\mathfrak{c}$}{valorc}
$\mathfrak{c} = |\mathbb{R}| = |\mathcal{P}(\mathbb{N})| = |2^{\mathbb{N}}| = 2^{\aleph_0}$, donde $2^{\mathbb{N}}$ denota el conjunto de todas las funciones $\mathbb{N} \to \{0,1\}$.
\end{proposicion}

\begin{proof}
Construimos biyecciones $\mathbb{R} \sim \mathcal{P}(\mathbb{N}) \sim \{0,1\}^{\mathbb{N}}$.

\textbf{$\{0,1\}^{\mathbb{N}} \sim \mathcal{P}(\mathbb{N})$:} La biyección es $f \mapsto f^{-1}(\{1\}) = \{n \in \mathbb{N} \mid f(n) = 1\}$, que identifica cada función $\{0,1\}^{\mathbb{N}}$ con su conjunto «soporte».

\textbf{$\mathcal{P}(\mathbb{N}) \preceq [0,1]$:} A cada $S \subseteq \mathbb{N}$ le asociamos el número $\sum_{n \in S} 3^{-(n+1)}$ (expansión en base 3 con dígitos 0 y 1). Esta función es inyectiva (de hecho, su imagen es el conjunto de Cantor ternario).

\textbf{$[0,1] \preceq \mathcal{P}(\mathbb{N})$:} Cada $x \in [0,1]$ tiene una expansión binaria $x = \sum_{n=0}^\infty b_n 2^{-(n+1)}$, $b_n \in \{0,1\}$. La función $x \mapsto \{n \mid b_n = 1\}$ es inyectiva (salvo los números diádicos con dos expansiones, pero se puede elegir siempre la que termina en 0).

Por CBS, $\mathcal{P}(\mathbb{N}) \sim [0,1] \sim \mathbb{R}$.
\end{proof}

% ===========================================================
\section{El Teorema de Cantor: estrictamente más grande}\label{sec:cantorpot}
% ===========================================================

\begin{teorema}{Teorema de Cantor}{cantorpot}
Para todo conjunto $A$, $|A| < |\mathcal{P}(A)|$. Es decir, el conjunto potencia siempre tiene cardinal estrictamente mayor que el conjunto original.
\end{teorema}

\begin{proof}
\textbf{Parte 1: $|A| \leq |\mathcal{P}(A)|$.} La función $a \mapsto \{a\}$ es una inyección de $A$ en $\mathcal{P}(A)$.

\textbf{Parte 2: $|A| \neq |\mathcal{P}(A)|$.} Supongamos que existe una biyección $f: A \to \mathcal{P}(A)$. Definimos el subconjunto \textbf{diagonal}:
\[
  D = \{a \in A \mid a \notin f(a)\}.
\]
Como $D \subseteq A$, tenemos $D \in \mathcal{P}(A)$. Como $f$ es sobreyectiva, existe $a_0 \in A$ con $f(a_0) = D$. Entonces:
\begin{itemize}[leftmargin=2em]
  \item Si $a_0 \in D$, entonces por definición de $D$: $a_0 \notin f(a_0) = D$. Contradicción.
  \item Si $a_0 \notin D$, entonces $a_0 \notin f(a_0) = D$, luego por la definición de $D$: $a_0 \in D$. Contradicción.
\end{itemize}
En cualquier caso se obtiene una contradicción. Por tanto, no puede existir tal biyección.
\end{proof}

\begin{corolario}{Jerarquía infinita de cardinales}{jerarquiacard}
Existe una sucesión estrictamente creciente de cardinales infinitos:
\[
  \aleph_0 \;=\; |\mathbb{N}| \;<\; |\mathcal{P}(\mathbb{N})| \;<\; |\mathcal{P}(\mathcal{P}(\mathbb{N}))| \;<\; \cdots
\]
En particular, no existe un «conjunto de todos los cardinales» ni un «infinito máximo».
\end{corolario}

\begin{ejemplo}{El teorema de Cantor para conjuntos finitos}{cantorfinito}
Para $A = \{1, 2, 3\}$, $|A| = 3$ y $|\mathcal{P}(A)| = 2^3 = 8 > 3$. El subconjunto diagonal en la prueba del Teorema~\ref{teo:cantorpot} depende de la función $f$ elegida. Por ejemplo, si $f(1) = \{2, 3\}$, $f(2) = \{1\}$, $f(3) = \{1, 2, 3\}$, entonces:
\begin{align*}
  D &= \{a \in A \mid a \notin f(a)\}.
\end{align*}
Verificamos cada elemento:
\begin{itemize}[leftmargin=2em]
  \item $1 \notin f(1) = \{2,3\}$: \textbf{verdadero}, luego $1 \in D$.
  \item $2 \notin f(2) = \{1\}$: \textbf{verdadero}, luego $2 \in D$.
  \item $3 \notin f(3) = \{1,2,3\}$: \textbf{falso} ($3 \in \{1,2,3\}$), luego $3 \notin D$.
\end{itemize}
Por tanto $D = \{1, 2\}$. El conjunto $D = \{1, 2\}$ no es ninguno de los valores $f(1), f(2), f(3)$, lo que confirma que $f$ no puede ser sobreyectiva.
\end{ejemplo}

\begin{observacion}{La Hipótesis del Continuo}{hipotesisCH}
El Teorema de Cantor establece $\aleph_0 < \mathfrak{c} = 2^{\aleph_0}$. La pregunta natural es: ¿existe algún cardinal $\kappa$ con $\aleph_0 < \kappa < \mathfrak{c}$? La respuesta es que esta pregunta —conocida como la \textbf{Hipótesis del Continuo} (HC): «no existe tal $\kappa$», i.e., $\mathfrak{c} = \aleph_1$— es \textbf{independiente de ZFC}: Gödel demostró en 1940 que HC es consistente con ZFC, y Cohen demostró en 1963 que su negación también lo es \cite{godel1940, cohen1966}. Este es uno de los resultados más profundos de la lógica matemática del siglo XX.
\end{observacion}

% ===========================================================
\section{Resumen y panorama}
% ===========================================================

\begin{importante}
\textbf{Resumen de los resultados principales de esta clase:}
\begin{enumerate}[leftmargin=2em, label=\textbf{\arabic*.}]
  \item Una función $f: A \to B$ es inyectiva si no colapsa elementos distintos; sobreyectiva si alcanza todo el codominio; biyectiva si es ambas cosas a la vez.
  \item La equipotencia ($A \sim B$: existe una biyección) es una relación de equivalencia.
  \item El Teorema de Cantor–Bernstein–Schröder: $|A| \leq |B|$ y $|B| \leq |A|$ implican $|A| = |B|$.
  \item $\mathbb{Z}$, $\mathbb{Q}$, $\mathbb{Z}[x]$ y $\mathbb{A}$ (los algebraicos) son todos numerables: $|\cdot| = \aleph_0$.
  \item $\mathbb{R}$ no es numerable (argumento diagonal de Cantor): $|\mathbb{R}| = \mathfrak{c} > \aleph_0$.
  \item Para todo $A$: $|A| < |\mathcal{P}(A)|$ (Teorema de Cantor), lo que genera una jerarquía infinita de cardinalidades.
  \item En particular: $\aleph_0 < \mathfrak{c} = 2^{\aleph_0}$.
\end{enumerate}
\end{importante}

La teoría de cardinalidades continuará en la Clase~5, donde introduciremos los números cardinales de von Neumann, la aritmética cardinal (suma, producto y potencia de cardinales) y exploraremos las consecuencias del Axioma de Elección sobre la comparabilidad de cardinales.

% ===========================================================
\section{Problemas y tareas}
% ===========================================================

\begin{enumerate}[leftmargin=2.5em, label=\textbf{\arabic*.}]

  \item \textbf{(Tipos de funciones.)}
  \begin{enumerate}[label=(\alph*)]
    \item Para cada una de las siguientes funciones $f: \mathbb{R} \to \mathbb{R}$, determine si es inyectiva, sobreyectiva, biyectiva, o ninguna de las anteriores. Justifique formalmente.
    \begin{enumerate}[label=(\roman*)]
      \item $f(x) = x^3$.
      \item $f(x) = x^2$.
      \item $f(x) = e^x$.
      \item $f(x) = \sin(x)$.
      \item $f(x) = x^3 - x$.
    \end{enumerate}
    \item Dé un ejemplo de una función $f: \mathbb{N} \to \mathbb{N}$ que sea sobreyectiva pero no inyectiva. ¿Puede ocurrir lo contrario (inyectiva pero no sobreyectiva)?
    \item Si $f: A \to B$ es biyectiva, demuestre directamente desde la definición (sin usar la Proposición~\ref{prop:compoinysupr}) que $f^{-1}: B \to A$ es biyectiva.
  \end{enumerate}

  \item \textbf{(Imágenes directas e inversas.)}
  \begin{enumerate}[label=(\alph*)]
    \item Sea $f: \mathbb{R} \to \mathbb{R}$, $f(x) = x^2$. Calcule:
    \begin{itemize}
      \item $f([-3, 2])$, $f^{-1}([1, 4])$, $f^{-1}(\{-1\})$.
    \end{itemize}
    \item Demuestre que para $f: A \to B$ y $S, S' \subseteq A$:
    \[
      f(S \cap S') = f(S) \cap f(S') \iff f \text{ es inyectiva}.
    \]
    \item Demuestre que $f(f^{-1}(T)) = T$ para todo $T \subseteq B$ si y solo si $f$ es sobreyectiva.
    \item (Desafío) Encuentre conjuntos $A, B$ y una función $f: A \to B$ tal que $f(A \setminus S) \neq f(A) \setminus f(S)$ para algún $S \subseteq A$.
  \end{enumerate}

  \item \textbf{(Dedekind-infinitud.)}
  \begin{enumerate}[label=(\alph*)]
    \item Demuestre que $\mathbb{Z}$ es Dedekind-infinito exhibiendo explícitamente una biyección entre $\mathbb{Z}$ y un subconjunto propio suyo.
    \item Demuestre que si $A$ es Dedekind-infinito y $A \subseteq B$, entonces $B$ es Dedekind-infinito.
    \item Sea $A$ un conjunto y $\iota: \mathbb{N} \to A$ una función inyectiva. Construya explícitamente una función inyectiva $f: A \to A$ que no sea sobreyectiva, demostrando que $A$ es Dedekind-infinito.
    \item (Investigación) Enuncie y explique por qué la equivalencia «infinito $\Leftrightarrow$ Dedekind-infinito» requiere el Axioma de Elección Numerable (una forma debilitada del Axioma de Elección). Consulte \cite{herrlich2006}.
  \end{enumerate}

  \item \textbf{(Equipotencia.)}
  \begin{enumerate}[label=(\alph*)]
    \item Demuestre directamente (sin CBS) que $[0,1] \sim [0,2]$ y que $[0,1] \sim [a,b]$ para cualquier $a < b$.
    \item Demuestre que $(0, +\infty) \sim \mathbb{R}$ construyendo una biyección explícita.
    \item Demuestre que $\mathbb{N} \sim \mathbb{N} \setminus \{0, 1, 2, \ldots, k\}$ para todo $k \in \mathbb{N}$.
    \item Sea $A = \{f: \mathbb{N} \to \{0,1,2\} \mid f \text{ es eventualmente } 0\}$ (funciones de $\mathbb{N}$ en $\{0,1,2\}$ que son $0$ salvo finitos valores). Determine si $A \sim \mathbb{N}$, $A \sim \mathbb{R}$, o ninguna de las dos. Justifique.
    \item Demuestre que $\mathbb{N}^{\mathbb{N}}$ (el conjunto de todas las sucesiones de naturales) cumple $|\mathbb{N}^{\mathbb{N}}| = |\mathbb{R}|$.
  \end{enumerate}

  \item \textbf{(Teorema de Cantor–Bernstein–Schröder.)}
  \begin{enumerate}[label=(\alph*)]
    \item Use CBS para demostrar que $[0,1] \sim [0,1)$ sin construir una biyección explícita. (Sugerencia: construya inyecciones en ambas direcciones.)
    \item Demuestre que el conjunto de los irracionales en $(0,1)$ es equipotente a $\mathbb{R}$.
    \item Sea $A \subseteq B \subseteq C$ con $A \sim C$. Demuestre que $A \sim B \sim C$. (Este es el «Teorema de Dedekind–Bernstein».)
    \item (Desafío) Demuestre que si $f: A \to B$ es sobreyectiva, entonces $|B| \leq |A|$. ¿Requiere AC este resultado?
  \end{enumerate}

  \item \textbf{(Numerabilidad.)}
  \begin{enumerate}[label=(\alph*)]
    \item Demuestre que el conjunto de los números racionales en $(0,1)$ es numerable, construyendo explícitamente una enumeración.
    \item Demuestre que el conjunto de todos los pares $(p, q)$ con $p, q \in \mathbb{Z}$ y $\gcd(p,q) = 1$ es numerable.
    \item Demuestre que el conjunto de todas las sucesiones finitas de $\{0,1\}$ (cadenas binarias finitas, incluyendo la cadena vacía) es numerable.
    \item Demuestre que el conjunto de todos los subconjuntos \emph{finitos} de $\mathbb{N}$ es numerable, pero el de todos los subconjuntos de $\mathbb{N}$ no lo es.
    \item Sea $A = \{r \in \mathbb{R} \mid r \text{ es solución de un polinomio de grado } \leq 5 \text{ con coeficientes en } \mathbb{Z}\}$. Determine si $A$ es numerable.
  \end{enumerate}

  \item \textbf{(Argumento diagonal.)}
  \begin{enumerate}[label=(\alph*)]
    \item Demuestre que el conjunto $\{0,1\}^{\mathbb{N}}$ de todas las sucesiones binarias infinitas no es numerable, usando el argumento diagonal directamente sobre sucesiones (sin hacer referencia a $\mathbb{R}$).
    \item Deduzca a partir del punto (a) que $(0,1)$ no es numerable, usando la representación binaria de los reales. ¿Qué hay que hacer con los números diádicos?
    \item (Paradoja de Berry) Considere la siguiente «definición»: «el menor número natural que no puede describirse en menos de veinte palabras en español». ¿Por qué esta definición es paradójica? Relaciones esto con el argumento diagonal.
    \item (Generalización) Demuestre el siguiente resultado: para cualquier conjunto $A$ y cualquier enumeración $f_0, f_1, f_2, \ldots$ de funciones $f_i: A \to \{0,1\}$, existe una función $g: A \to \{0,1\}$ que no aparece en la lista. ¿Qué condición se necesita sobre $A$?
  \end{enumerate}

  \item \textbf{(Teorema de Cantor: $|A| < |\mathcal{P}(A)|$.)}
  \begin{enumerate}[label=(\alph*)]
    \item Demuestre el Teorema de Cantor para $A = \mathbb{N}$: construya explícitamente un subconjunto $D \subseteq \mathbb{N}$ que no sea la imagen de ningún elemento bajo una función hipotética $f: \mathbb{N} \to \mathcal{P}(\mathbb{N})$ sobreyectiva.
    \item Demuestre que no existe ninguna función sobreyectiva de $A$ en $\mathcal{P}(A)$, sin suponer que $f$ sea inyectiva. (El teorema de Cantor solo requiere la ausencia de sobreyección, no la ausencia de inyección.)
    \item Use el Teorema de Cantor para probar que no existe un «conjunto universal» en ZFC, es decir, no puede existir un conjunto $V$ tal que $A \in V$ para todo conjunto $A$.
    \item Deduzca que la clase propia $\mathbf{V}$ de todos los conjuntos no puede ser un conjunto.
  \end{enumerate}

  \item \textbf{(Proyecto: la Hipótesis del Continuo.)}
  Investigue y redacte un ensayo breve (2--3 páginas) sobre los siguientes puntos:
  \begin{itemize}[leftmargin=2em]
    \item Enuncie la Hipótesis del Continuo: $\nexists \kappa$ con $\aleph_0 < \kappa < \mathfrak{c}$, equivalentemente $\mathfrak{c} = \aleph_1$.
    \item Explique en qué sentido Gödel (1940) demostró que HC es \emph{consistente} con ZFC, usando el constructo de los «conjuntos construibles» $\mathbf{L}$.
    \item Explique en qué sentido Cohen (1963) demostró que la negación de HC también es consistente con ZFC, usando la técnica del \emph{forcing}.
    \item ¿Qué significa que HC sea \emph{independiente} de ZFC? ¿Es esto un defecto de ZFC?
    \item ¿Cuál es el estado actual de la investigación sobre la verdad de HC?
  \end{itemize}
  \textbf{Fuentes sugeridas:} \cite{cohen1966, jech2003, kunen2011, kanamori2003}.

\end{enumerate}

% ===========================================================
\section*{Referencias bibliográficas de esta clase}
% ===========================================================

Los resultados desarrollados en esta clase se encuentran tratados con rigor y detalle en las siguientes fuentes:

\begin{itemize}[leftmargin=2em]
  \item \textbf{Tratamientos estándar de funciones, equipotencia y cardinales:} \cite{enderton1977} ofrece una presentación muy clara y rigurosa de funciones, equipotencia y los primeros cardinales transfinitos, accesible al nivel de este curso; \cite{halmos1960} es un clásico indispensable; \cite{hrbacek1999} es un texto introductorio completo y recomendado.

  \item \textbf{Teoría de conjuntos avanzada:} \cite{jech2003} es la referencia canónica para todos los aspectos de la teoría de cardinales y ordinales; \cite{kunen2011} es preciso y exhaustivo en los aspectos formales.

  \item \textbf{Teorema de Cantor–Bernstein–Schröder:} El resultado aparece en su forma moderna en \cite{enderton1977} y \cite{jech2003}; la demostración vía el Teorema del Punto Fijo de Knaster–Tarski se puede encontrar en \cite{devlin1993} y en muchos textos de álgebra de retículos.

  \item \textbf{Argumento diagonal y no numerabilidad de $\mathbb{R}$:} El artículo original de Cantor es \cite{cantor1891}; una presentación moderna y accesible se encuentra en \cite{enderton1977} y \cite{sierpinski1958}.

  \item \textbf{Conjuntos Dedekind-infinitos y Axioma de Elección:} \cite{herrlich2006} es una monografía dedicada íntegramente a las consecuencias y patologías del Axioma de Elección; \cite{moore1982} ofrece el contexto histórico.

  \item \textbf{Hipótesis del Continuo:} \cite{cohen1966} es el libro clásico de Cohen donde se presenta el forcing; \cite{godel1940} es el trabajo original de Gödel; \cite{kanamori2003} sitúa estos resultados en el panorama más amplio de los cardinales grandes.

  \item \textbf{Definición de Dedekind de conjunto infinito:} \cite{dedekind1888} es la fuente histórica original donde Dedekind define los conjuntos infinitos mediante la reflexividad.
\end{itemize}
